\begin{proof}[Proof of \cref{obj:lem:projection-control}]
Write \(\|\cdot\|_2\) and \(\langle\cdot,\cdot\rangle\) for the norm and inner product in \(L^2(\lambda)\). We use the cosine basis and projection kernel of \cref{obj:synth_3}.

\begin{enumerate}
\item Fix \(0<\gamma\le1\), a continuous function \(f:[0,1]\to\mathbb R\), and an integer \(k\ge1\). If \([f]_\gamma=+\infty\), then
\[
5[f]_\gamma k^{-\gamma}=+\infty,
\]
because \(k^{-\gamma}>0\), and the extended-real approximation inequality follows. Otherwise define
\[
L_f=[f]_\gamma\in[0,\infty).
\]
The definition of the seminorm gives
\[
|f(x)-f(z)|\le L_f|x-z|^\gamma
\qquad(x,z\in[0,1]).
\]
For \(x\ne z\), this follows by multiplying the defining quotient bound by the positive denominator; for \(x=z\), both sides vanish.
% lean: Causalean.Mathlib.Analysis.Approximation.Chebyshev.CosineProjection.cosineProjection_error_le_seminorm

\item Define the folding map and the even extension by
\[
\chi(\vartheta)=\frac{\arccos(\cos\vartheta)}{\pi},
\qquad
\bar f(\vartheta)=f(\chi(\vartheta)),
\qquad \vartheta\in\mathbb R.
\]
The map \(\chi\) takes values in \([0,1]\), is continuous, even, and \(2\pi\)-periodic, and satisfies
\[
\chi(\pi x)=x\qquad(x\in[0,1]).
\]
We also have
\[
|\chi(\vartheta)-\chi(\vartheta')|
\le \frac{|\vartheta-\vartheta'|}{\pi}
\qquad(\vartheta,\vartheta'\in\mathbb R).
\]
To see this, choose an integer \(\ell_{\mathrm{per}}\) such that
\[
\vartheta'-2\pi\ell_{\mathrm{per}}\in[-\pi,\pi).
\]
The elementary identities and bounds
\[
\arccos(\cos\vartheta)=|\vartheta|
\quad\text{if }|\vartheta|\le\pi,
\qquad
\arccos(\cos\vartheta)\le|\vartheta|
\quad\text{for every }\vartheta\in\mathbb R
\]
give, by periodicity and the triangle inequality,
\[
\begin{aligned}
\arccos(\cos\vartheta)
&\le|\vartheta-2\pi\ell_{\mathrm{per}}|\\
&\le|\vartheta'-2\pi\ell_{\mathrm{per}}|
   +|\vartheta-\vartheta'|\\
&=\arccos(\cos\vartheta')+|\vartheta-\vartheta'|.
\end{aligned}
\]
Interchanging the two angles proves the asserted Lipschitz bound. Consequently,
\[
|\bar f(\vartheta)-\bar f(\vartheta')|
\le L_f\left(\frac{|\vartheta-\vartheta'|}{\pi}\right)^\gamma.
\]
The extension \(\bar f\) is continuous, even, and \(2\pi\)-periodic.
% lean: Causalean.Mathlib.Analysis.Approximation.Chebyshev.CosineProjection.evenExtension_holder

\item For an integer \(m_{\mathrm J}\ge1\), introduce the continuous sine quotient
\[
\mathcal S_{m_{\mathrm J}}(\vartheta)
=
\begin{cases}
\dfrac{\sin(m_{\mathrm J}\vartheta/2)}{\sin(\vartheta/2)},
   &\sin(\vartheta/2)\ne0,\\[6pt]
m_{\mathrm J}(-1)^{(m_{\mathrm J}-1)\iota_{\mathrm{per}}},
   &\vartheta=2\pi\iota_{\mathrm{per}},\quad \iota_{\mathrm{per}}\in\mathbb Z,
\end{cases}
\qquad \vartheta\in\mathbb R,
\]
and its mass and normalized fourth-power kernel
\[
\mathcal Z_{m_{\mathrm J}}
=\int_{-\pi}^{\pi}\mathcal S_{m_{\mathrm J}}(\vartheta)^4\,d\vartheta,
\qquad
\mathcal J_{m_{\mathrm J}}(\vartheta)
=\frac{\mathcal S_{m_{\mathrm J}}(\vartheta)^4}{\mathcal Z_{m_{\mathrm J}}}.
\]
The removable values in the definition are the corresponding limits.

The finite cosine identity is
\[
\mathcal S_{m_{\mathrm J}}(\vartheta)^2
=m_{\mathrm J}
 +2\sum_{j=1}^{m_{\mathrm J}-1}(m_{\mathrm J}-j)\cos(j\vartheta).
\]
For completeness, the Chebyshev polynomials of the second kind, defined by
\[
\mathsf U_0(z)=1,\qquad
\mathsf U_1(z)=2z,\qquad
\mathsf U_{j+1}(z)=2z\mathsf U_j(z)-\mathsf U_{j-1}(z)
\quad(j\ge1),
\]
satisfy
\[
\mathcal S_{m_{\mathrm J}}(\vartheta)
=\mathsf U_{m_{\mathrm J}-1}(\cos(\vartheta/2)).
\]
We first derive the polynomial identity
\[
\mathsf U_j(z)^2-\mathsf U_{j-1}(z)^2=\mathsf U_{2j}(z)
\qquad(j\ge1).
\]
Regard these real-coefficient polynomials also as polynomials on \(\mathbb C\). Surjectivity of the complex cosine allows us, for each complex argument, to choose a complex angle satisfying
\[
z_{\mathrm C}=\cos\vartheta_{\mathrm C},
\qquad z_{\mathrm C},\vartheta_{\mathrm C}\in\mathbb C.
\]
If \(\sin\vartheta_{\mathrm C}=0\), then \(\cos\vartheta_{\mathrm C}=1\) or \(-1\). Induction in the defining recurrence gives
\[
\mathsf U_j(1)=j+1,
\qquad
\mathsf U_j(-1)=(-1)^j(j+1)
\qquad(j\ge0).
\]
Thus, at either argument,
\[
\mathsf U_j(\pm1)^2-\mathsf U_{j-1}(\pm1)^2
=(j+1)^2-j^2=2j+1=\mathsf U_{2j}(\pm1)
\qquad(j\ge1).
\]
If \(\sin\vartheta_{\mathrm C}\ne0\), induction in the same recurrence, using the sine addition formula, gives
\[
\mathsf U_j(\cos\vartheta_{\mathrm C})
=\frac{\sin((j+1)\vartheta_{\mathrm C})}
       {\sin\vartheta_{\mathrm C}}
\qquad(j\ge0).
\]
Expanding the sine and cosine addition formulas and using
\(\cos^2\vartheta_{\mathrm C}+\sin^2\vartheta_{\mathrm C}=1\) yields
\[
\sin((j+1)\vartheta_{\mathrm C})^2
-\sin(j\vartheta_{\mathrm C})^2
=\sin((2j+1)\vartheta_{\mathrm C})
 \sin\vartheta_{\mathrm C}.
\]
Dividing by \(\sin^2\vartheta_{\mathrm C}\) proves
\[
\begin{aligned}
\mathsf U_j(\cos\vartheta_{\mathrm C})^2
-\mathsf U_{j-1}(\cos\vartheta_{\mathrm C})^2
&=\frac{\sin((2j+1)\vartheta_{\mathrm C})}
        {\sin\vartheta_{\mathrm C}}\\
&=\mathsf U_{2j}(\cos\vartheta_{\mathrm C}).
\end{aligned}
\]
The two cases establish the polynomial identity at every complex argument, hence also over the reals.

For the even-index cosine expansion, introduce the first-kind Chebyshev polynomials by
\[
\mathsf T^{\mathrm{Ch}}_0(z_{\mathrm{Ch}})=1,
\qquad
\mathsf T^{\mathrm{Ch}}_1(z_{\mathrm{Ch}})=z_{\mathrm{Ch}},
\qquad z_{\mathrm{Ch}}\in\mathbb R,
\]
\[
\mathsf T^{\mathrm{Ch}}_{j+1}(z_{\mathrm{Ch}})
=2z_{\mathrm{Ch}}\mathsf T^{\mathrm{Ch}}_j(z_{\mathrm{Ch}})
 -\mathsf T^{\mathrm{Ch}}_{j-1}(z_{\mathrm{Ch}})
\qquad(j\ge1).
\]
The cosine addition formula and induction give
\[
\mathsf T^{\mathrm{Ch}}_j(\cos(\vartheta/2))
=\cos(j\vartheta/2)
\qquad(j\ge0,\ \vartheta\in\mathbb R).
\]
The defining recurrences also imply
\[
\mathsf U_{j+2}(z_{\mathrm{Ch}})-\mathsf U_j(z_{\mathrm{Ch}})
=2\mathsf T^{\mathrm{Ch}}_{j+2}(z_{\mathrm{Ch}})
\qquad(j\ge0).
\]
Indeed, both sides satisfy the same second-order recurrence as sequences in \(j\), and their first two values agree:
\[
\begin{aligned}
\mathsf U_2(z_{\mathrm{Ch}})-\mathsf U_0(z_{\mathrm{Ch}})
&=4z_{\mathrm{Ch}}^2-2
 =2\mathsf T^{\mathrm{Ch}}_2(z_{\mathrm{Ch}}),\\
\mathsf U_3(z_{\mathrm{Ch}})-\mathsf U_1(z_{\mathrm{Ch}})
&=8z_{\mathrm{Ch}}^3-6z_{\mathrm{Ch}}
 =2\mathsf T^{\mathrm{Ch}}_3(z_{\mathrm{Ch}}).
\end{aligned}
\]
Consequently,
\[
\mathsf U_{2j+2}(\cos(\vartheta/2))
=\mathsf U_{2j}(\cos(\vartheta/2))
 +2\cos((j+1)\vartheta)
\qquad(j\ge0).
\]
The base case is \(\mathsf U_0=1\). If the expansion holds at \(j\), the last identity gives
\[
\begin{aligned}
\mathsf U_{2j+2}(\cos(\vartheta/2))
&=1+2\sum_{\iota_{\mathrm{cos}}=1}^{j}
 \cos(\iota_{\mathrm{cos}}\vartheta)
 +2\cos((j+1)\vartheta)\\
&=1+2\sum_{\iota_{\mathrm{cos}}=1}^{j+1}
 \cos(\iota_{\mathrm{cos}}\vartheta).
\end{aligned}
\]
Thus induction proves
\[
\mathsf U_{2j}(\cos(\vartheta/2))
=1+2\sum_{\iota_{\mathrm{cos}}=1}^{j}
 \cos(\iota_{\mathrm{cos}}\vartheta)
\qquad(j\ge0),
\]
with an empty sum at \(j=0\).
Starting with \(\mathsf U_0^2=1\) and summing these differences yields the displayed cosine identity.

Integrating that identity, and then its square, gives
\[
\int_{-\pi}^{\pi}\mathcal S_{m_{\mathrm J}}(\vartheta)^2\,d\vartheta
=2\pi m_{\mathrm J},
\]
and
\[
\begin{aligned}
\mathcal Z_{m_{\mathrm J}}
&=2\pi m_{\mathrm J}^2
  +4\pi\sum_{j=1}^{m_{\mathrm J}-1}(m_{\mathrm J}-j)^2\\
&=2\pi m_{\mathrm J}^2
  +\frac{4\pi}{6}(m_{\mathrm J}-1)m_{\mathrm J}(2m_{\mathrm J}-1)\\
&=\frac{2\pi}{3}m_{\mathrm J}(2m_{\mathrm J}^2+1)
 \ge\frac{4\pi}{3}m_{\mathrm J}^3>0.
\end{aligned}
\]
Here nonconstant cosines integrate to zero, distinct cosine modes are orthogonal, and each positive-frequency squared cosine has integral \(\pi\). In particular,
\[
\mathcal J_{m_{\mathrm J}}(\vartheta)\ge0,
\qquad
\int_{-\pi}^{\pi}\mathcal J_{m_{\mathrm J}}(\vartheta)\,d\vartheta=1.
\]
These formulas include \(m_{\mathrm J}=1\), when the sums are empty.
% lean: Causalean.Mathlib.Analysis.JacksonApproximation.jrawMass_eq

\item For \(|\vartheta|\le\pi\), the sine chord bound gives
\[
|\vartheta|\le\pi|\sin(\vartheta/2)|.
\]
Together with
\[
\mathcal S_{m_{\mathrm J}}(\vartheta)\sin(\vartheta/2)
=\sin(m_{\mathrm J}\vartheta/2),
\]
this implies
\[
\vartheta^2\mathcal S_{m_{\mathrm J}}(\vartheta)^2
\le\pi^2\sin(m_{\mathrm J}\vartheta/2)^2
\le\pi^2.
\]
The identity also holds at the removable points. Multiplying by
\(\mathcal S_{m_{\mathrm J}}(\vartheta)^2\) and integrating yields
\[
\int_{-\pi}^{\pi}
 \vartheta^2\mathcal S_{m_{\mathrm J}}(\vartheta)^4\,d\vartheta
\le
\pi^2\int_{-\pi}^{\pi}
 \mathcal S_{m_{\mathrm J}}(\vartheta)^2\,d\vartheta
=2\pi^3m_{\mathrm J}.
\]
Using the mass lower bound from the preceding step, we obtain
\[
\begin{aligned}
\int_{-\pi}^{\pi}\vartheta^2\mathcal J_{m_{\mathrm J}}(\vartheta)\,d\vartheta
&\le \frac{2\pi^3m_{\mathrm J}}{\mathcal Z_{m_{\mathrm J}}}\\
&\le \frac{2\pi^3m_{\mathrm J}}
              {(4\pi/3)m_{\mathrm J}^3}
\le \left(\frac{2\pi}{m_{\mathrm J}}\right)^2.
\end{aligned}
\]
The kernel therefore defines a probability measure
\[
d\nu_{\mathrm J}(\vartheta)
=\mathbf 1_{[-\pi,\pi]}(\vartheta)
 \mathcal J_{m_{\mathrm J}}(\vartheta)\,d\vartheta.
\]
Jensen's inequality for the concave function \(z\mapsto z^{\gamma/2}\) on \([0,\infty)\) gives
\[
\begin{aligned}
\int_{-\pi}^{\pi}
 \left(\frac{|\vartheta|}{\pi}\right)^\gamma
 \mathcal J_{m_{\mathrm J}}(\vartheta)\,d\vartheta
&=
\int_{\mathbb R}
 \left(\frac{\vartheta^2}{\pi^2}\right)^{\gamma/2}
 \,d\nu_{\mathrm J}(\vartheta)\\
&\le
\left(\int_{\mathbb R}\frac{\vartheta^2}{\pi^2}
 \,d\nu_{\mathrm J}(\vartheta)\right)^{\gamma/2}\\
&\le\left(\frac{2}{m_{\mathrm J}}\right)^\gamma.
\end{aligned}
\]
This argument applies throughout \(0<\gamma\le1\).
% lean: Causalean.Mathlib.Analysis.Approximation.Chebyshev.CosineProjection.jackson_holder_moment

\item Define the Jackson approximant on the unit interval by
\[
p_{\mathrm J}(x)
=\int_{-\pi}^{\pi}
 \bar f(\pi x-\vartheta)\mathcal J_{m_{\mathrm J}}(\vartheta)\,d\vartheta,
\qquad x\in[0,1].
\]
Normalization of the kernel and the increment bound for \(\bar f\) imply
\[
\begin{aligned}
|f(x)-p_{\mathrm J}(x)|
&=\left|\int_{-\pi}^{\pi}
  \{f(x)-\bar f(\pi x-\vartheta)\}
  \mathcal J_{m_{\mathrm J}}(\vartheta)\,d\vartheta\right|\\
&\le L_f\int_{-\pi}^{\pi}
 \left(\frac{|\vartheta|}{\pi}\right)^\gamma
 \mathcal J_{m_{\mathrm J}}(\vartheta)\,d\vartheta\\
&\le L_f\left(\frac{2}{m_{\mathrm J}}\right)^\gamma.
\end{aligned}
\]
The approximant is continuous.

It also lies in the cosine span through frequency \(2m_{\mathrm J}-2\). Indeed,
\(\mathsf U_{m_{\mathrm J}-1}^4\) is an even polynomial of degree \(4m_{\mathrm J}-4\); hence there is a polynomial \(\mathcal Q_{m_{\mathrm J}}\) such that
\[
\mathsf U_{m_{\mathrm J}-1}(z)^4
=\mathcal Q_{m_{\mathrm J}}(z^2),
\qquad
\deg\mathcal Q_{m_{\mathrm J}}\le2m_{\mathrm J}-2.
\]
Since \(\cos(\vartheta/2)^2=(1+\cos\vartheta)/2\),
\[
\mathcal J_{m_{\mathrm J}}(\vartheta)
=
\frac{\mathcal Q_{m_{\mathrm J}}((1+\cos\vartheta)/2)}
     {\mathcal Z_{m_{\mathrm J}}}.
\]
Thus the kernel is an even trigonometric polynomial. Averaging its trigonometric expansion at \(\vartheta\) and \(-\vartheta\) removes the sine terms, giving real coefficients
\(\mathfrak a_{m_{\mathrm J},j}\), \(0\le j\le2m_{\mathrm J}-2\), with
\[
\mathcal J_{m_{\mathrm J}}(\vartheta)
=\sum_{j=0}^{2m_{\mathrm J}-2}
 \mathfrak a_{m_{\mathrm J},j}\cos(j\vartheta).
\]
Reversing the convolution over one period and using the cosine subtraction formula gives
\[
\begin{aligned}
p_{\mathrm J}(x)
&=\int_{-\pi}^{\pi}
 \bar f(\vartheta)\mathcal J_{m_{\mathrm J}}(\pi x-\vartheta)\,d\vartheta\\
&=\sum_{j=0}^{2m_{\mathrm J}-2}
 \left\{\mathfrak a_{m_{\mathrm J},j}
 \int_{-\pi}^{\pi}\bar f(\vartheta)\cos(j\vartheta)\,d\vartheta\right\}
 \cos(\pi jx).
\end{aligned}
\]
The sine integrals vanish because \(\bar f\) is even. The reversal follows by a change of variable and periodicity. Since \(\varphi_0=1\) and
\(\varphi_j(x)=\sqrt2\cos(\pi jx)\) for \(j\ge1\), this is the required cosine-span representation.
% lean: Causalean.Mathlib.Analysis.Approximation.Chebyshev.CosineProjection.jacksonApproximant_mem_span

\item Now take
\[
m_{\mathrm J}=\left\lfloor\frac{k+1}{2}\right\rfloor.
\]
For every \(k\ge1\),
\[
m_{\mathrm J}\ge1,\qquad
2m_{\mathrm J}-2<k,\qquad
k\le2m_{\mathrm J}.
\]
Consequently \(p_{\mathrm J}\) belongs to the range of \(\Pi_k\), and
\[
\left(\frac{2}{m_{\mathrm J}}\right)^\gamma
\le\left(\frac4k\right)^\gamma
=4^\gamma k^{-\gamma}
\le5k^{-\gamma},
\]
where \(4^\gamma\le4\le5\). Hence
\[
|f(x)-p_{\mathrm J}(x)|\le5L_fk^{-\gamma}
\qquad(x\in[0,1]).
\]
Orthogonality of the projection gives
\[
\|f-p_{\mathrm J}\|_2^2
=\|f-\Pi_kf\|_2^2+\|\Pi_kf-p_{\mathrm J}\|_2^2,
\]
so, because \(\lambda\) has total mass one,
\[
\|f-\Pi_kf\|_2
\le\|f-p_{\mathrm J}\|_2
\le5L_fk^{-\gamma}.
\]
For \(k=1\), the chosen order is \(m_{\mathrm J}=1\) and the approximant is constant, so the same argument applies. In the finite-seminorm case these ordinary norms equal the extended norms in the statement; together with the infinite-seminorm case, this proves the approximation assertion.
% lean: CausalSmith.Stat.LogoddsLowsmoothFrontier.cosineProjection_error_le_holderSeminorm

\item Fix measurable marks \(W,V\) valued in \([0,1]\) and \(n\ge2\). Let \(O,O'\) be independent records and let \(O_1,\ldots,O_n\) be the original independent sample:
\[
(O,O')\sim P\otimes P,
\qquad
(O_1,\ldots,O_n)\sim P^{\otimes n}.
\]
Writing \(p_{ay}(x)=P(A=a,Y=y\mid X=x)\) for the specified four-cell probabilities, the conditional means are
\[
w(x)=\sum_{a,y\in\{0,1\}}p_{ay}(x)W(x,a,y),
\qquad
v(x)=\sum_{a,y\in\{0,1\}}p_{ay}(x)V(x,a,y).
\]
Their measurability and the normalization of the cells give
\[
0\le w(x),v(x)\le1,\qquad x\in[0,1],
\]
so both belong to \(L^2(\lambda)\). Disintegration and
\cref{obj:ass:uniform-design} give, for every integrable measurable \(h:[0,1]\to\mathbb R\),
\[
E_P[W(O)h(X)]=\int_0^1w(x)h(x)\,d\lambda(x),
\qquad
E_P[V(O)h(X)]=\int_0^1v(x)h(x)\,d\lambda(x).
\]

Define the symmetric record kernel
\[
q_k(o,o')
=\frac{H_k(x,x')}{2}
 \{W(o)V(o')+V(o)W(o')\},
\qquad x=X(o),\quad x'=X(o').
\]
Swapping the two indices in the complete ordered-pair sum shows that
\[
\mathbb U_{n,k}(W,V)
=\frac1{n(n-1)}\sum_{i\ne j}q_k(O_i,O_j).
\]
The finite kernel expansion and the preceding conditional-mean identities yield
\[
\int q_k(o,o')\,dP(o')
=\frac{W(o)(\Pi_kv)(X(o))+V(o)(\Pi_kw)(X(o))}{2}.
\]
Indeed,
\[
\begin{aligned}
\int H_k(x,X(o'))V(o')\,dP(o')
&=\sum_{j=0}^{k-1}\varphi_j(x)
  \int_0^1v(z)\varphi_j(z)\,d\lambda(z)\\
&=(\Pi_kv)(x),
\end{aligned}
\]
and the corresponding identity holds with \(W,w\).
% lean: CausalSmith.Stat.LogoddsLowsmoothFrontier.symmetricProjectionKernel_integral_right

\item Each distinct pair of sample coordinates has law \(P\otimes P\), and there are \(n(n-1)\) such ordered pairs. Therefore
\[
E_P\mathbb U_{n,k}(W,V)=E_Pq_k(O,O').
\]
Integrating the kernel section from the preceding step gives
\[
\begin{aligned}
E_Pq_k(O,O')
&=\frac12\{\langle w,\Pi_kv\rangle+\langle v,\Pi_kw\rangle\}\\
&=\sum_{j=0}^{k-1}
 \langle w,\varphi_j\rangle\langle v,\varphi_j\rangle\\
&=\langle\Pi_kw,\Pi_kv\rangle.
\end{aligned}
\]
Here the last equality uses the orthonormality supplied by
\cref{obj:synth_3}. All sums are finite; bounded marks and the continuous finite kernel ensure integrability.
% lean: CausalSmith.Stat.LogoddsLowsmoothFrontier.projectionStatistic_integral

\item The same finite expansion gives
\[
\langle w,\Pi_kv\rangle
=\langle v,\Pi_kw\rangle
=\langle\Pi_kw,\Pi_kv\rangle.
\]
Expanding the product of the residuals thus yields
\[
\begin{aligned}
\langle w-\Pi_kw,v-\Pi_kv\rangle
&=\langle w,v\rangle-\langle w,\Pi_kv\rangle
 -\langle v,\Pi_kw\rangle+\langle\Pi_kw,\Pi_kv\rangle\\
&=\langle w,v\rangle-\langle\Pi_kw,\Pi_kv\rangle\\
&=\langle w,v\rangle-E_P\mathbb U_{n,k}(W,V).
\end{aligned}
\]
This proves the bias identity.
% lean: CausalSmith.Stat.LogoddsLowsmoothFrontier.projectionStatistic_bias

\item For the variance argument define, for all records \(o,o'\),
\[
\begin{aligned}
\tau_k&=\langle\Pi_kw,\Pi_kv\rangle,\\
\psi_k(o)&=\int q_k(o,o')\,dP(o'),\\
\mathcal R_k(o,o')
&=q_k(o,o')-\psi_k(o)-\psi_k(o')+\tau_k.
\end{aligned}
\]
Also define the two finite-sample terms
\[
\mathbb A_{n,k}
=\frac2n\sum_{i=1}^n\{\psi_k(O_i)-\tau_k\},
\qquad
\mathbb B_{n,k}
=\frac1{n(n-1)}\sum_{i\ne j}\mathcal R_k(O_i,O_j).
\]
These quantities are square integrable: the finite cosine kernel is bounded on the compact covariate square, the marks are bounded, and all sums are finite.

The kernel mean identity gives \(E_P\psi_k(O)=\tau_k\). Hence
\[
\int\mathcal R_k(o,o')\,dP(o')
=\psi_k(o)-\psi_k(o)-\tau_k+\tau_k=0.
\]
Symmetry of \(q_k\) gives symmetry of \(\mathcal R_k\), so
\[
\int\mathcal R_k(o,o')\,dP(o)=0,
\qquad
E_P\mathcal R_k(O,O')=0.
\]
For \(i\ne j\) and \(1\le\iota_{\mathrm{sample}}\le n\),
\[
E_P\!\left[
 \{\psi_k(O_{\iota_{\mathrm{sample}}})-\tau_k\}\mathcal R_k(O_i,O_j)
\right]=0.
\]
If \(\iota_{\mathrm{sample}}=i\), integrate first over \(O_j\); if \(\iota_{\mathrm{sample}}=j\), integrate first over \(O_i\). In each case the corresponding residual section integrates to zero. If \(\iota_{\mathrm{sample}}\) differs from both indices, independence factors the expectation, and the residual mean is zero. Expanding the finite product of sums therefore gives
\[
E_P[\mathbb A_{n,k}\mathbb B_{n,k}]=0.
\]
% lean: CausalSmith.Stat.LogoddsLowsmoothFrontier.projection_hoeffding_cross

\item The explicit formula for the first projection implies
\[
\psi_k(o)^2
\le\frac{(\Pi_kv)(X(o))^2+(\Pi_kw)(X(o))^2}{2}.
\]
Indeed,
\[
\begin{aligned}
&\{W(o)(\Pi_kv)(X(o))+V(o)(\Pi_kw)(X(o))\}^2\\
&\qquad\le2\{W(o)^2(\Pi_kv)(X(o))^2
             +V(o)^2(\Pi_kw)(X(o))^2\},
\end{aligned}
\]
by the elementary square inequality, and the squared marks are at most one. Applying the residual identity with the same conditional mean in both positions gives
\[
\|\Pi_kw\|_2^2
=\|w\|_2^2-\|w-\Pi_kw\|_2^2
\le\|w\|_2^2\le1,
\]
and likewise \(\|\Pi_kv\|_2^2\le1\). Uniform design consequently yields
\[
E_P\psi_k(O)^2
\le\frac{\|\Pi_kv\|_2^2+\|\Pi_kw\|_2^2}{2}\le1.
\]
Expanding the centered square and using \(E_P\psi_k(O)=\tau_k\), we obtain
\[
E_P\{\psi_k(O)-\tau_k\}^2
=E_P\psi_k(O)^2-\tau_k^2\le1.
\]
Independence makes the cross products between distinct centered sample terms vanish. Thus
\[
\begin{aligned}
E_P\mathbb A_{n,k}^2
&=\frac4{n^2}\,
 nE_P\{\psi_k(O)-\tau_k\}^2\\
&=\frac4nE_P\{\psi_k(O)-\tau_k\}^2
\le\frac4n.
\end{aligned}
\]
% lean: CausalSmith.Stat.LogoddsLowsmoothFrontier.projectionFirst_finite_secondMoment_le

\item The mark restrictions give the pointwise inequality
\[
q_k(o,o')^2\le H_k(X(o),X(o'))^2.
\]
By the finite cosine expansion and orthonormality,
\[
\int_0^1 H_k(x,z)^2\,d\lambda(z)
=\sum_{j=0}^{k-1}\varphi_j(x)^2,
\]
and hence
\[
E_Pq_k(O,O')^2
\le\int_0^1\!\int_0^1H_k(x,z)^2\,d\lambda(z)\,d\lambda(x)
=\sum_{j=0}^{k-1}\|\varphi_j\|_2^2=k.
\]
To obtain the residual energy, note that integration of a kernel section gives
\[
E_P[q_k(O,O')\psi_k(O)]
=E_P[q_k(O,O')\psi_k(O')]
=E_P\psi_k(O)^2.
\]
Independence also gives
\[
E_P[\psi_k(O)\psi_k(O')]=\tau_k^2,
\qquad
E_Pq_k(O,O')=\tau_k.
\]
Expanding the defining square of \(\mathcal R_k\) and substituting these identities yields
\[
E_P\mathcal R_k(O,O')^2
=E_Pq_k(O,O')^2-2E_P\psi_k(O)^2+\tau_k^2.
\]
Using the centered-energy identity from the preceding step, this becomes
\[
\begin{aligned}
E_P\mathcal R_k(O,O')^2
&=E_Pq_k(O,O')^2-E_P\psi_k(O)^2
  -E_P\{\psi_k(O)-\tau_k\}^2\\
&\le E_Pq_k(O,O')^2\le k.
\end{aligned}
\]
% lean: CausalSmith.Stat.LogoddsLowsmoothFrontier.projectionResidual_energy_le_rank

\item For two ordered distinct pairs \((i,j)\) and \((i',j')\), residual degeneracy gives
\[
E_P[
 \mathcal R_k(O_i,O_j)\mathcal R_k(O_{i'},O_{j'})
]
=
\begin{cases}
E_P\mathcal R_k(O,O')^2,
 &\{i,j\}=\{i',j'\},\\
0,&\{i,j\}\ne\{i',j'\}.
\end{cases}
\]
For equal index sets, the pairs are identical or reversed, and symmetry gives the first value. For unequal index sets, one record appears in just one factor; integrating over that record first gives zero. This covers both disjoint pairs and pairs sharing one record.

Each of the \(n(n-1)\) ordered pairs has exactly two ordered pairs with the same index set. Expanding the square of the sum therefore gives
\[
E_P\left[
 \left\{\sum_{i\ne j}\mathcal R_k(O_i,O_j)\right\}^2
\right]
=2n(n-1)E_P\mathcal R_k(O,O')^2.
\]
Equivalently,
\[
nE_P\mathbb B_{n,k}^2
=\frac2{n-1}E_P\mathcal R_k(O,O')^2.
\]
Since \(n\ge2\), division by \(n\) and the residual energy bound yield
\[
E_P\mathbb B_{n,k}^2
=\frac{2E_P\mathcal R_k(O,O')^2}{n(n-1)}
\le\frac{2k}{n(n-1)}.
\]
The counting includes \(n=2\), when the two ordered pairs are reversals of one another.
% lean: CausalSmith.Stat.LogoddsLowsmoothFrontier.projectionResidual_finite_secondMoment_le

\item Finally, summing the defining residual identity over ordered distinct pairs gives
\[
\sum_{i\ne j}\mathcal R_k(O_i,O_j)
=\sum_{i\ne j}q_k(O_i,O_j)
 -2(n-1)\sum_{i=1}^n\psi_k(O_i)+n(n-1)\tau_k.
\]
Each first-coordinate term and each second-coordinate term occurs \(n-1\) times. Dividing by \(n(n-1)\) proves the pointwise decomposition
\[
\mathbb U_{n,k}(W,V)-\tau_k
=\mathbb A_{n,k}+\mathbb B_{n,k}.
\]
Since \(E_P\mathbb U_{n,k}(W,V)=\tau_k\), expansion of the square, the zero cross moment, and the two established bounds give
\[
\begin{aligned}
\operatorname{Var}_P\{\mathbb U_{n,k}(W,V)\}
&=E_P(\mathbb A_{n,k}+\mathbb B_{n,k})^2\\
&=E_P\mathbb A_{n,k}^2+E_P\mathbb B_{n,k}^2
  +2E_P[\mathbb A_{n,k}\mathbb B_{n,k}]\\
&\le\frac4n+\frac{2k}{n(n-1)}.
\end{aligned}
\]
Together with the approximation, mean, and bias identities, this proves the result.
% lean: CausalSmith.Stat.LogoddsLowsmoothFrontier.projectionStatistic_variance_le
\end{enumerate}
\end{proof}
